% Figure 6: Catastrophic Failure Detection (K=5 Portfolio)
% Uses production router with semantic transfer for all 5 models
% Updated: February 2026

\section{Experiment: Corralling as Safety Mechanism}
\label{sec:corralling}

\subsection{Motivation: Automatic Failover as a Secondary Benefit}

Production LLM systems face a critical challenge: \textbf{models can catastrophically fail with minimal warning}. API endpoints crash, providers silently update models with quality regressions, or traffic shifts to domains where warmup priors perform poorly.

\textbf{The Question}: Does Corralling---designed primarily for safe expert hedging and contextual routing---provide useful automatic failover when a model catastrophically fails? And how does this scale with portfolio size?

\subsection{Experimental Setup: 5-Model Portfolio}

\subsubsection{Why $K=5$?}

With $K=2$ models, catastrophic failure detection is a trivial binary tracking problem: a simple EMA tracker just needs to learn ``model A bad, model B good'' from per-model averages. With $K=5$, post-failure routing becomes \emph{contextual}: the optimal policy must redistribute traffic across 4 remaining models based on prompt features. Each non-failing model excels in a different region of context space, so per-model averages cannot distinguish which alternative is best for a given prompt.

\subsubsection{Production Router Components}

\begin{itemize}
    \item \textbf{CostAwareLinUCBRouter} (warmup expert): Initialized with real RouteLLM priors for Mixtral and GPT-4-Turbo; semantic transfer (First-Child Bias Correction) for GPT-3.5-Turbo, Claude-3-Haiku, and GPT-4o
    \item \textbf{CostAwareTabulaRasaRouter} (tabula rasa expert): Cold-start LinUCB with ridge regularization for all 5 models
    \item \textbf{CorrallingRouter}: Coordinates experts with importance-weighted meta-learning via selection tokens
\end{itemize}

All 5 models have informed warmup priors. Both experts explore freely across all models.

\subsubsection{Three-Phase Scenario}

Context-dependent reward environment (orthogonal context directions per model, noise $\sigma=0.08$):

\begin{itemize}
    \item \textbf{Phase 1 ($t=0$--$100$):} All 5 models healthy
    \item \textbf{Phase 2 ($t=100$--$300$):} GPT-4-Turbo catastrophically fails ($\mu=0.15$, Cohen's $d\approx 5.0$)
    \item \textbf{Phase 3 ($t=300$--$500$):} GPT-4-Turbo recovers
\end{itemize}

Corralling settings: $\eta=0.3$, $\gamma=0.05$. Baselines: Oracle, Static GPT-4-Turbo, EMA Tracker ($\alpha=0.15$, $\epsilon=0.1$).

\subsection{Results ($N=20$ Seeds)}

\begin{figure}[t]
\centering
\includegraphics[width=0.95\linewidth]{results/figure6_5model.pdf}
\caption{\textbf{Catastrophic Failure Detection: 5-Model Portfolio} ($K=5$, $N=20$ seeds).
\textbf{(Top)} Running-average reward. EMA reacts faster ($\sim 0.74$) than banditGPT ($\sim 0.65$) but concentrates 63\% of traffic on Mixtral alone, missing context-dependent alternatives.
\textbf{(Middle)} Expert weight evolution: 95\% detection rate with clear tabula rasa dominance during failure.
\textbf{(Bottom)} Model selection fractions reveal EMA's $\epsilon/K = 2\%$ per-model exploration limit at $K=5$.}
\label{fig:corralling}
\end{figure}

\begin{table}[h]
\centering
\caption{Mean reward by phase ($K=5$, $N=20$ seeds)}
\begin{tabular}{lccc}
\toprule
\textbf{Method} & \textbf{Healthy} & \textbf{Failure} & \textbf{Recovery} \\
\midrule
Oracle & 0.824 & 0.800 & 0.822 \\
\textbf{banditGPT} & \textbf{0.763} & \textbf{0.646} & \textbf{0.762} \\
EMA Tracker & 0.736 & 0.735 & 0.783 \\
Static GPT-4-Turbo & 0.812 & 0.153 & 0.820 \\
\bottomrule
\end{tabular}
\end{table}

\subsection{Portfolio Size Scaling}

\begin{table}[h]
\centering
\caption{Gap narrows monotonically with portfolio size}
\begin{tabular}{lcccc}
\toprule
$K$ & banditGPT & EMA & $\Delta$ & Detection \\
\midrule
2 & 0.478 & 0.764 & $-0.286$ & 65\% \\
3 & 0.596 & 0.750 & $-0.154$ & 85\% \\
\textbf{5} & \textbf{0.646} & \textbf{0.735} & $\mathbf{-0.090}$ & \textbf{95\%} \\
\bottomrule
\end{tabular}
\end{table}

The gap reduces by 69\% from $K=2$ to $K=5$. EMA degrades as $K$ grows (exploration budget $\epsilon/K$ thins across more models); Corralling's contextual experts become more valuable as the routing decision becomes higher-dimensional.

\subsection{Interpretation}

Catastrophic failure detection is a \textbf{secondary safety benefit} of the Corralling mechanism, not its primary design objective. For pure model-level failure detection, dedicated monitoring (e.g., EMA tracking, CUSUM) should complement Corralling. However, the monotonic gap closure with portfolio size suggests that at production-scale portfolios ($K \geq 8$), Corralling's contextual advantage may fully compensate for its detection latency.

\subsection{Reproducibility}

Code: \texttt{experiments\_v1/06\_figure/generate\_figure6\_5model.py}
\begin{itemize}
    \item Uses production router with semantic transfer for all 5 models
    \item Configuration: $\eta=0.3$, $\gamma=0.05$, $N_{\text{eff}}=10$, 500 steps, 20 seeds
    \item Runtime: $\sim$7 seconds
\end{itemize}

\bibliographystyle{ACM-Reference-Format}
\begin{thebibliography}{1}

\bibitem{agarwal2017corralling}
Alekh Agarwal, Haipeng Luo, Behnam Neyshabur, and Robert E. Schapire.
\newblock Corralling a band of bandit algorithms.
\newblock In \emph{Conference on Learning Theory (COLT)}, pages 12--38. PMLR, 2017.

\end{thebibliography}
